\documentclass[11pt,letterpaper]{article}
\usepackage[utf8]{inputenc}
\usepackage[T1]{fontenc}
\usepackage{lmodern}
\usepackage[spanish,es-noshorthands,es-tabla]{babel}
\usepackage{csquotes}
\usepackage[backend=biber,style=numeric]{biblatex}
\addbibresource{referencias.bib}
\usepackage{amsmath}
\usepackage{amssymb}
\usepackage{graphicx}
\usepackage{fancyhdr}
\usepackage[hyperindex=true,breaklinks=true,colorlinks=true,linkcolor=black,urlcolor=blue,citecolor=black]{hyperref}
\usepackage{booktabs}
\usepackage{tabularx}
\usepackage{array}
\usepackage{anysize}
\usepackage{caption}
\usepackage[table]{xcolor}
\usepackage{float}
\usepackage{listings}

\newcommand{\SafeIncludeGraphics}[2][]{\IfFileExists{#2}{\includegraphics[#1]{#2}}{\mbox{}}}
\newcommand{\code}[1]{\texttt{#1}}

% --- Listados de codigo ---------------------------------------------------
\definecolor{comentario}{gray}{0.40}
\lstdefinelanguage{rvasm}{
  sensitive=false,
  morecomment=[l]{\#},
  morekeywords={add,addi,sub,and,andi,or,ori,xor,xori,slt,slti,sltu,sltiu,
    sll,slli,srl,srli,sra,srai,lui,auipc,jal,jalr,beq,bne,blt,bge,bltu,bgeu,
    lb,lh,lw,lbu,lhu,sb,sh,sw,li,mv,j,call,ret,beqz,bnez,bleu,nop},
}
\lstset{
  basicstyle=\ttfamily\footnotesize,
  keywordstyle=\bfseries,
  commentstyle=\color{comentario},
  columns=fullflexible,
  keepspaces=true,
  frame=single,
  breaklines=true,
  showstringspaces=false,
  captionpos=b,
}
\renewcommand{\lstlistingname}{Código}
\renewcommand{\lstlistlistingname}{Índice de códigos}

%% Encabezado y pie ------------------------------------------------------------
\pagestyle{fancy}
\fancyhf{}
\chead{Arquitectura de Computadores\vspace{0.3cm}}
\lhead{\SafeIncludeGraphics[scale=0.4]{uandes.pdf}}
\lfoot{\scriptsize Facultad Ingenier\'ia y Ciencias Aplicadas\\Universidad de los Andes}
\rfoot{\thepage}
\renewcommand{\footrulewidth}{0.25pt}
\addtolength{\headheight}{1.4cm}
\renewcommand{\baselinestretch}{1.3}

\begin{document}
%% Portada ----------------------------------------------------------------------
\begin{titlepage}
\begin{center}
\SafeIncludeGraphics[width=1.2in]{Log.png}\\[0.8cm]
\large UNIVERSIDAD DE LOS ANDES \\
\large Facultad de Ingeniería y Ciencias Aplicadas \\
Arquitectura de Computadores 2026-20\\
\vspace{2cm}
{\LARGE \textbf{Proyecto Nº2}}\\
\vspace{0.4cm}
{\large Juego de reflejos en assembly sobre Pochoco SoC}\\
\vspace{1.5cm}
Grupo 2:\\
Josefa Alcaíno Lira \\
Felipe Izcúe Ferrer \\
Martín Norton Haigis\\
\vspace{1.5cm}
Profesor de Cátedra: Jorge Gómez M. \\
\vspace{0.1cm}
Ayudantes: \\ Domingo Fuenzalida \\
\vspace{0.1cm} Matthias Ruhe\\
\vspace{0.5cm}
29 de septiembre de 2026
\end{center}
\end{titlepage}

\marginsize{2.5cm}{2cm}{0.8cm}{2cm}
\newpage
\begin{abstract}
Se implementó un juego de reflejos en assembly RV32E que corre sobre el procesador
Espino del Pochoco SoC, instanciado en una FPGA Lattice iCE40 HX1K de la placa Nandland
Go Board. En cada ronda el juego enciende los cuatro LEDs durante tres segundos, elige
un LED objetivo de forma pseudoaleatoria, mide el tiempo que tarda el jugador en
apretar el botón correspondiente y lo muestra en décimas de segundo en los dos displays
de siete segmentos. Tras diez respuestas correctas muestra el promedio.

Para medir tiempo se agregó al SoC un contador de ciclos de 32 bits, mapeado en memoria
en la dirección \code{0x8000000C}; el cambio ocupa 96 celdas lógicas y el diseño
completo usa el 90\,\% de la FPGA, con una frecuencia máxima de 46~MHz frente a los
25~MHz del reloj. Como el procesador no tiene multiplicación, división ni
desplazamientos, las divisiones se hacen por restas sucesivas y el generador
pseudoaleatorio, un LFSR de 32 bits, se construyó solo con sumas y comparaciones.

El programa se traduce a código de máquina con un assembler propio escrito en Python,
que implementa RV32I completo con los registros de RV32E, pseudo-instrucciones,
constantes con nombre y validación de cada operando. El assembler se verificó con
20 pruebas automáticas, entre ellas la reproducción exacta de la salida del assembler
oficial de GNU y la comprobación de todos los inmediatos posibles contra las fórmulas
del decodificador del procesador. El juego se verificó en simulación sobre el RTL
completo del SoC, con un testbench que juega dos partidas y revisa cada regla del
enunciado mirando solo los pines de la placa.
\end{abstract}

\newpage
\tableofcontents
\newpage

\section{Introducción}

El objetivo del proyecto es programar en assembly un juego de reflejos que corra sobre
un procesador RISC-V real, instanciado en una FPGA, y traducir ese programa a código de
máquina con un assembler desarrollado por el grupo. El punto de partida es el Pochoco
SoC \cite{pochoco}: un sistema en chip con el procesador Espino (RV32E), una memoria RAM
unificada para instrucciones y datos, y periféricos mapeados en memoria para los LEDs,
los botones y los displays de la Go Board \cite{goboard}.

El trabajo tiene tres partes, que corresponden a las secciones de este informe:

\begin{enumerate}
  \item \textbf{Hardware} (sección~\ref{sec:arquitectura}). El SoC original no ofrece
        ninguna forma de medir tiempo. Se agregó un contador de ciclos de reloj que el
        programa puede leer, y un módulo superior \code{game\_top} que instancia el SoC
        con la imagen de memoria del juego.
  \item \textbf{Software} (secciones~\ref{sec:temporizacion} y~\ref{sec:programa}). El
        juego en \code{sw/game.s}: temporización de tres segundos, medición del tiempo
        de respuesta, conversión de ciclos a décimas de segundo, promedio de las diez
        rondas y elección pseudoaleatoria del LED objetivo.
  \item \textbf{Assembler} (sección~\ref{sec:assembler}). El programa
        \code{asm/pochoasm.py}, que traduce \code{game.s} a \code{game.hex}: una palabra
        de 32 bits en hexadecimal por línea, que la RAM del SoC carga al sintetizar el
        diseño.
\end{enumerate}

La sección~\ref{sec:verificacion} describe cómo se verificó cada parte, y la
sección~\ref{sec:operacion} cómo compilar, programar la placa y cambiar una constante
del juego, que es lo que se pedirá durante la demostración.

\subsection{Reglas del juego}
\label{sec:reglas}

La secuencia implementada sigue el enunciado:

\begin{enumerate}
  \item Al encender, los LEDs parpadean alternando 0101 y 1010 y los displays muestran
        \code{00}. El juego espera que se apriete cualquier botón para empezar.
  \item Cada ronda comienza con los cuatro LEDs encendidos durante al menos tres
        segundos.
  \item Luego se enciende solo el LED objetivo, elegido de forma pseudoaleatoria. En
        ese mismo instante comienza la medición de tiempo.
  \item El jugador aprieta un botón. Si es el del LED objetivo, los displays muestran
        el tiempo de respuesta en décimas de segundo y la ronda cuenta como correcta.
  \item Si es otro botón, los displays muestran \code{EE} y la ronda se repite desde
        el paso~2. Las rondas correctas acumuladas se conservan.
  \item Tras diez respuestas correctas, el último tiempo queda visible tres segundos
        y después se muestra el promedio, con los LEDs parpadeando. Un botón inicia
        una partida nueva.
\end{enumerate}

Sobre el paso~3, el enunciado pide iniciar la medición ``en el mismo instante lógico en
que comienza la ronda''. Se interpretó que el tiempo de respuesta se cuenta desde que
aparece el estímulo, es decir, desde que se enciende el LED objetivo: si se contara
desde que se encienden los cuatro LEDs, todos los tiempos incluirían los tres segundos
fijos de la primera fase. La lectura del contador que inicia la medición es la
instrucción inmediatamente siguiente a la que enciende el LED objetivo.

\section{Arquitectura del sistema}
\label{sec:arquitectura}

\subsection{Pochoco SoC y procesador Espino}

El procesador Espino implementa RV32E \cite{riscv-unpriv}: el conjunto de instrucciones
base de 32 bits de RISC-V con solo 16 registros (\code{x0} a \code{x15}). No incluye
la extensión M (multiplicación y división). Además, el ALU del repositorio tiene los
desplazamientos deshabilitados para ahorrar celdas lógicas: \code{sll}, \code{srl},
\code{sra} y sus variantes con inmediato se decodifican bien, pero se ejecutan como una
suma. Estas dos restricciones determinan cómo se escribió el programa
(sección~\ref{sec:sin-mul}).

El core tiene dos etapas: búsqueda de instrucción (IF) y decodificación-ejecución (ID).
El banco de registros está hecho con bloques de RAM de la FPGA, que entregan el dato un
ciclo después de pedirlo, así que cada instrucción espera un ciclo por sus operandos.
Se midió en simulación la duración de cada tipo de instrucción, leyendo el contador de
ciclos antes y después de ella (tabla~\ref{tab:cpi}). Estos valores se usan en la
sección~\ref{sec:temporizacion} para acotar el error de la medición.

\begin{table}[H]
\centering
\caption{Ciclos de reloj por instrucción del Espino, medidos en simulación.}
\label{tab:cpi}
\begin{tabular}{lc}
\toprule
Instrucción & Ciclos \\
\midrule
ALU (\code{add}, \code{addi}, \code{xor}, \ldots), \code{lui}, \code{sw} & 2 \\
Branch no tomado & 2 \\
\code{lw} (RAM o periférico) & 3 \\
Branch tomado, \code{jal}, \code{jalr} & 3 \\
\bottomrule
\end{tabular}
\end{table}

La RAM tiene 512 palabras (2\,KB) y guarda a la vez el programa y los datos. Se
inicializa al sintetizar con el contenido de un archivo \code{.hex} mediante
\code{\$readmemh}, y el procesador comienza a ejecutar en la dirección 0. Los
periféricos se seleccionan con el bit 31 de la dirección
(tabla~\ref{tab:mapa}).

\begin{table}[H]
\centering
\caption{Mapa de memoria usado por el juego. La última fila es el registro agregado.}
\label{tab:mapa}
\begin{tabular}{lll}
\toprule
Dirección & Acceso & Contenido \\
\midrule
\code{0x0000\_0000} & lectura/escritura & RAM: programa y datos (512 palabras) \\
\code{0x8000\_0000} & escritura & displays: bits 7:4 izquierdo, bits 3:0 derecho \\
\code{0x8000\_0004} & escritura & LEDs, bits 3:0 \\
\code{0x8000\_0008} & lectura & botones, bits 3:0 (1 = apretado) \\
\code{0x8000\_000C} & lectura & \textbf{contador de ciclos, 32 bits} \\
\bottomrule
\end{tabular}
\end{table}

Cada display muestra un dígito hexadecimal: el periférico traduce los 4 bits a los
siete segmentos. Para mostrar un número decimal de dos cifras el programa escribe
$\text{decenas}\cdot 16 + \text{unidades}$; por ejemplo, para 37 escribe \code{0x37}.

\subsection{Modificaciones al RTL}

Se hicieron tres cambios al hardware, todos visibles con \code{git diff} desde el primer
commit del repositorio, que es el código original sin modificar. El tercero corrige un
defecto del SoC original que solo se manifestó en la placa
(sección~\ref{sec:placa}).

\paragraph{Contador de ciclos.} El juego necesita medir 3 segundos y tiempos de reacción.
Un bucle de espera calibrado no sirve para medir, porque su duración depende de qué
instrucciones ejecuta (tabla~\ref{tab:cpi}) y cambiaría con cualquier modificación al
programa. Se agregó a \code{rtl/pochoco\_periph.v} un registro de 32 bits que suma uno
en cada flanco de reloj y que el programa lee en el desplazamiento \code{0x0C} del
bloque de periféricos (código~\ref{lst:contador}). Es de solo lectura, comienza en cero
tras el reset y da la vuelta cada $2^{32}$ ciclos, unos 171,8 segundos.

\begin{lstlisting}[language=Verilog, caption={Contador de ciclos agregado a \code{pochoco\_periph.v}.}, label={lst:contador}]
  reg [31:0] cycle_q;

  always @(posedge clk_i or negedge rst_ni) begin
    if (!rst_ni) cycle_q <= 32'b0;
    else         cycle_q <= cycle_q + 32'd1;
  end
  ...
      case (off)
        6'd2: rdata_o <= {28'b0, btn_i}; // Buttons
        6'd3: rdata_o <= cycle_q;        // Cycle counter   <- nuevo
\end{lstlisting}

\paragraph{Módulo superior \code{game\_top}.} El archivo \code{rtl/game\_top.v} es el
que se sintetiza. Tiene los mismos puertos que la Go Board según \code{goboard.pcf} y
solo instancia el SoC con el parámetro \code{MemFile} apuntando a \code{sw/game.hex}.
El botón $i$ corresponde al LED $i$: \code{i\_Switch[i]} y \code{o\_LED[i]}.

\paragraph{Reset de encendido.} En \code{rtl/pochoco\_soc.v}, el contador que mantiene
al procesador en reset después de configurar la FPGA pasó de 4 a 12 bits: el procesador
arranca a los 4.095 ciclos (164~$\mu$s) en vez de a los 15 (0,6~$\mu$s). Con el valor
original, el juego no funcionaba en la placa aunque todas las simulaciones pasaban. La
sección~\ref{sec:placa} describe cómo se encontró la causa.

\begin{lstlisting}[language=Verilog, caption={Reset de encendido en \code{pochoco\_soc.v}, con el contador ampliado.}, label={lst:por}]
  reg  [11:0] por_cnt = 12'b0;           // antes: reg [3:0], hasta 4'hF
  wire        rst_ni;
  always @(posedge clk) begin
    if (por_cnt != 12'hFFF) por_cnt <= por_cnt + 12'd1;
  end
  assign rst_ni = (por_cnt == 12'hFFF);
\end{lstlisting}

\paragraph{Lo que no se cambió.} Se evaluó activar los desplazamientos en el ALU, pero
el SoC original ya ocupaba el 83\,\% de las celdas lógicas de la FPGA, y un
desplazador de 32 bits cuesta del orden de cien celdas. El juego se escribió sin
desplazamientos (sección~\ref{sec:sin-mul}).

\subsection{Recursos y frecuencia}

La tabla~\ref{tab:recursos} compara el SoC original con el diseño final, ambos
sintetizados con Yosys y nextpnr \cite{osscad}. El contador de ciclos agrega 96 celdas
lógicas: 32 flip-flops, la cadena de suma y la entrada nueva del multiplexor de lectura.
Los 8 bits extra del reset de encendido agregan otras 14. La frecuencia máxima que reporta
nextpnr sigue muy por encima de los 25~MHz del reloj de la Go Board, y el analizador
independiente \code{icetime} de IceStorm da un resultado equivalente. El Makefile ejecuta nextpnr con \code{--freq 25}, así que la compilación falla
si alguna vez no se alcanzara esa frecuencia.

\begin{table}[H]
\centering
\caption{Uso de la FPGA iCE40 HX1K.}
\label{tab:recursos}
\begin{tabular}{lcc}
\toprule
 & SoC original & Diseño final \\
\midrule
Celdas lógicas (de 1280) & 1063 (83\,\%) & 1173 (91\,\%) \\
Bloques de RAM (de 16) & 12 & 12 \\
Frecuencia máxima & 45,0~MHz & 42,8~MHz \\
\bottomrule
\end{tabular}
\end{table}

\section{Temporización y medición de tiempo}
\label{sec:temporizacion}

\subsection{Frecuencia de reloj y conversión de ciclos a décimas}

La Go Board tiene un oscilador de 25~MHz conectado al pin 15 de la FPGA, y el SoC usa
ese reloj sin dividirlo. Todas las cantidades de tiempo del juego se derivan de ese
valor:

\begin{align*}
  f_{\text{clk}} &= 25\,000\,000~\text{Hz}, \qquad T_{\text{clk}} = \frac{1}{f_{\text{clk}}} = 40~\text{ns} \\[4pt]
  N_{\text{décima}} &= f_{\text{clk}} \cdot 0{,}1~\text{s} = 2\,500\,000~\text{ciclos} \\[4pt]
  N_{3\,\text{s}} &= f_{\text{clk}} \cdot 3~\text{s} = 75\,000\,000~\text{ciclos}
\end{align*}

Si el contador marca $c_0$ cuando se enciende el LED objetivo y $c_1$ cuando se detecta
el botón, el tiempo de respuesta en décimas de segundo es

\[
  t = \left\lfloor \frac{(c_1 - c_0) \bmod 2^{32}}{2\,500\,000} \right\rfloor .
\]

\begin{itemize}
  \item \textbf{La resta se hace sin signo, en 32 bits} (\code{sub}, que da el
        resultado módulo $2^{32}$). Si el contador dio la vuelta entre $c_0$ y $c_1$,
        la resta sigue dando la cantidad correcta de ciclos. El único caso que falla
        es una espera de más de $2^{32}$ ciclos, unos 171,8 segundos.
  \item \textbf{La división trunca}, como un cronómetro: 0,29~s se muestra como
        \code{02}. Así un valor mostrado de $n$ décimas significa que el tiempo real
        está entre $n$ y $n+1$ décimas.
  \item \textbf{La división se hace por restas sucesivas} (subrutina \code{divu}),
        porque el procesador no tiene instrucción de división. Cada vuelta del bucle
        toma 9 ciclos, así que convertir 9,9~s (99 vueltas) toma unos 900 ciclos,
        36~$\mu$s. Esa cuenta ocurre después de leer $c_1$, así que no afecta la
        medición.
\end{itemize}

Las constantes están con nombre al comienzo de \code{game.s}
(código~\ref{lst:constantes}). Cambiar la frecuencia del reloj solo exige cambiar
\code{WAIT\_CYCLES} y \code{TENTH\_CYCLES}.

\begin{lstlisting}[language=rvasm, caption={Constantes de tiempo en \code{sw/game.s}.}, label={lst:constantes}]
.equ ROUNDS,          10          # respuestas correctas por partida
.equ WAIT_CYCLES,     75000000    # fase con los 4 LEDs: 3 s = 3 * 25_000_000
.equ TENTH_CYCLES,    2500000     # ciclos en una decima de segundo = 25_000_000 / 10
.equ DEBOUNCE_CYCLES, 500000      # 20 ms con los botones sueltos antes de mostrar el objetivo
.equ BLINK_CYCLES,    6250000     # 0,25 s: cada cuanto cambian los LEDs mientras se espera
\end{lstlisting}

\subsection{Espera de tres segundos}

La subrutina \code{wait\_cycles} (código~\ref{lst:wait}) lee el contador al entrar y
vuelve cuando la diferencia con la lectura actual alcanza el valor pedido en
\code{a0}. La comparación \code{bltu} es sin signo, por la misma razón que la resta.

\begin{lstlisting}[language=rvasm, caption={Espera de \code{a0} ciclos.}, label={lst:wait}]
wait_cycles:
    lw   t0, CYCLES(s0)             # t0 = inicio
wc_loop:
    lw   t1, CYCLES(s0)
    sub  t1, t1, t0                 # t1 = ciclos transcurridos
    bltu t1, a0, wc_loop
    ret
\end{lstlisting}

Cada vuelta del bucle toma 8 ciclos (\code{lw} 3, \code{sub} 2, \code{bltu} tomado 3),
así que la espera termina a lo más 8 ciclos, 320~ns, después de cumplirse los
75.000.000. Tras ella el programa exige además 20~ms seguidos sin ningún botón
apretado (siguiente subsección). La fase con los cuatro LEDs dura entonces
$3{,}02$~s si nadie toca los botones, y más si alguien los mantiene apretados, lo que
cumple el ``al menos tres segundos'' del enunciado.

\subsection{Rebotes de los botones}

Un botón mecánico rebota: al apretarlo o soltarlo, el contacto se abre y cierra varias
veces durante algunos milisegundos. El juego solo necesita el primer contacto para
detener el cronómetro, así que los rebotes al apretar no afectan la medición. Los
rebotes al soltar sí son un problema. Si el jugador suelta el botón justo cuando
aparece el LED objetivo, un rebote se leería como una pulsación nueva, con un tiempo de
casi cero.

La subrutina \code{wait\_release} se ejecuta antes de encender el LED objetivo. Vuelve
solo cuando los botones llevan \code{DEBOUNCE\_CYCLES} (20~ms) seguidos sin apretarse:
cada lectura con un botón apretado reinicia la cuenta. Así se filtran los rebotes y,
además, no se puede ``hacer trampa'' dejando un botón apretado antes de que aparezca el
objetivo.

\subsection{Medición del tiempo de respuesta}

El código~\ref{lst:medicion} muestra el núcleo de la ronda. El cronómetro se inicia
con la instrucción siguiente a la que enciende el LED objetivo, y se detiene con la
instrucción siguiente a la que detecta el botón.

\begin{lstlisting}[language=rvasm, caption={Encendido del objetivo y medición.}, label={lst:medicion}]
    sw   a3, LEDS(s0)               #    solo el objetivo encendido
    lw   a2, CYCLES(s0)             # 3) y en ese mismo instante parte el cronometro
wait_press:
    lw   t0, BUTTONS(s0)
    beqz t0, wait_press             # 4) esperar un boton
    lw   t1, CYCLES(s0)             #    fin del cronometro
    bne  t0, a3, wrong              #    boton incorrecto (o mas de uno a la vez)
    sub  a0, t1, a2                 # 5) ciclos de reaccion
\end{lstlisting}

El error de la medición se acota con la tabla~\ref{tab:cpi}:

\begin{itemize}
  \item Entre que se escribe el LED y se lee $c_0$ pasan del orden de 2 ciclos.
  \item El bucle \code{wait\_press} da una vuelta cada 6 ciclos (\code{lw} 3 y
        \code{beqz} tomado 3), así que el botón se detecta a lo más 6 ciclos después de
        apretarse.
  \item Entre la detección y la lectura de $c_1$ pasan unos 4 ciclos.
\end{itemize}

En total el error es de unos 10 ciclos, 0,4~$\mu$s, doscientas mil veces menor que la
resolución de una décima. La medición es, en la práctica, exacta al ciclo que la
precede.

\subsection{Promedio}
\label{sec:promedio}

El programa acumula en el registro \code{a4} las décimas de cada ronda correcta. Tras
\code{ROUNDS} respuestas correctas, el promedio se calcula con la misma subrutina de
división:

\[
  \bar{t} = \left\lfloor \frac{\sum_{i=1}^{\text{ROUNDS}} t_i}{\text{ROUNDS}} \right\rfloor~\text{décimas}.
\]

Se promedian los tiempos en décimas y no los ciclos porque así el promedio es el de los
diez números que el jugador vio en los displays. Además, la suma de diez tiempos en
ciclos puede desbordar 32 bits si las rondas son lentas: $2^{32}$ ciclos son 171,8~s en
total. La suma en décimas no tiene ese problema.

Los displays tienen dos dígitos, así que un tiempo de 10~s o más se muestra como
\code{99}. Esa saturación es solo de la visualización: a la suma entra el valor
verdadero. Por ejemplo, en la simulación de la sección~\ref{sec:verificacion}, una
ronda de 12,0~s se muestra como \code{99} pero suma 120 décimas al promedio.

Las rondas con botón incorrecto no modifican ni \code{a4} ni el contador de rondas
\code{a5}: la ronda se repite y solo el intento correcto aporta su tiempo.

\section{El programa en assembly}
\label{sec:programa}

El juego completo está en \code{sw/game.s}, reproducido en el
anexo~\ref{anx:game}. Ocupa 108 líneas de instrucciones, que el assembler expande a
114 palabras de máquina (456 de los 2048 bytes de la RAM).

\subsection{Estructura}

El programa se divide en un cuerpo principal con dos bucles anidados, partida y ronda,
y seis subrutinas (tabla~\ref{tab:subrutinas}).

\begin{enumerate}
  \item \code{\_start}: carga la base de periféricos en \code{s0} y pone los displays
        en \code{00}.
  \item \code{game}: espera que se suelten los botones, parpadea hasta que se apriete
        uno (\code{wait\_start}), toma ese instante como semilla del generador
        pseudoaleatorio y reinicia el contador de rondas (\code{a5}) y la suma de
        tiempos (\code{a4}).
  \item \code{round}: enciende los cuatro LEDs, espera \code{WAIT\_CYCLES}, espera
        botones sueltos, elige el objetivo (\code{pick\_target}), lo enciende y lee el
        contador.
  \item \code{wait\_press}: espera un botón y lee el contador otra vez. Si el botón es
        incorrecto salta a \code{wrong}, que escribe \code{0xEE} en los displays y
        vuelve a \code{round}. Si es correcto, convierte los ciclos a décimas, suma,
        muestra el tiempo e incrementa \code{a5}.
  \item Si \code{a5} $<$ \code{ROUNDS} vuelve a \code{round}. Si no, apaga los LEDs,
        deja ver el último tiempo durante \code{WAIT\_CYCLES}, muestra el promedio y
        vuelve a \code{game}, donde \code{wait\_start} lo mantiene en pantalla mientras
        los LEDs parpadean.
\end{enumerate}

\begin{table}[H]
\centering
\caption{Subrutinas de \code{game.s}. Todas son hoja: no llaman a otras.}
\label{tab:subrutinas}
\begin{tabularx}{\linewidth}{l l X}
\toprule
Subrutina & Entrada / salida & Qué hace \\
\midrule
\code{wait\_cycles} & \code{a0} = ciclos & espera \code{a0} ciclos leyendo el contador \\
\code{wait\_release} & -- & vuelve cuando los botones llevan \code{DEBOUNCE\_CYCLES} sin apretarse \\
\code{wait\_start} & sale \code{a0} = contador & parpadea 0101/1010 hasta que se apriete un botón \\
\code{pick\_target} & sale \code{a0} = LED one-hot & avanza el LFSR (\code{s1}) 32 pasos y elige un LED \\
\code{divu} & \code{a0}, \code{a1} $\to$ \code{a0} & división entera sin signo por restas sucesivas \\
\code{show\_number} & \code{a0} = número & muestra \code{a0} en decimal en los displays, máximo 99 \\
\bottomrule
\end{tabularx}
\end{table}

\subsection{Uso de registros}

RV32E tiene solo 16 registros, así que se asignaron de forma fija
(tabla~\ref{tab:registros}). Como ninguna subrutina llama a otra, el registro de
retorno \code{ra} nunca se sobrescribe antes de un \code{ret}, y el programa no
necesita stack: \code{sp}, \code{gp} y \code{tp} quedan sin usar.

\begin{table}[H]
\centering
\caption{Asignación de registros.}
\label{tab:registros}
\begin{tabular}{lll}
\toprule
Registro & Contenido & Vida \\
\midrule
\code{s0} & base de periféricos \code{0x80000000} & todo el programa \\
\code{s1} & estado del LFSR & una partida \\
\code{a2} & contador al encender el objetivo ($c_0$) & una ronda \\
\code{a3} & LED objetivo en one-hot & una ronda \\
\code{a4} & suma de décimas de las rondas correctas & una partida \\
\code{a5} & rondas correctas & una partida \\
\code{a0}, \code{a1} & argumentos y resultado de subrutinas & una llamada \\
\code{t0}--\code{t2} & temporales & una llamada \\
\bottomrule
\end{tabular}
\end{table}

\subsection{Aritmética sin multiplicación, división ni desplazamientos}
\label{sec:sin-mul}

\begin{itemize}
  \item \textbf{Multiplicar por 2} es sumar un número consigo mismo:
        \code{add x, x, x}. Para escribir las decenas en los 4 bits altos del display
        se multiplican por 16 con cuatro sumas seguidas.
  \item \textbf{Desplazar un bit a la izquierda} es lo mismo que multiplicar por 2. Así
        avanza el LFSR (sección~\ref{sec:lfsr}) y así se construye la máscara del LED
        objetivo: se parte de \code{1} y se duplica tantas veces como indique el índice.
  \item \textbf{Dividir} es contar cuántas veces cabe el divisor restándolo
        (\code{divu}). Es lento para cocientes grandes, pero los cocientes del juego son
        pequeños: 99 como máximo para un tiempo y 99 para el promedio.
  \item \textbf{Separar decenas y unidades} para los displays es otra división por
        restas: se resta 10 mientras se pueda, y lo que queda son las unidades.
  \item \textbf{Leer el bit 31} de un registro no requiere desplazarlo: el bit 31 es el
        bit de signo, así que vale 1 exactamente cuando el registro es negativo
        (\code{bge s1, zero, ...}).
\end{itemize}

\subsection{Generación pseudoaleatoria}
\label{sec:lfsr}

El LED objetivo se elige con un LFSR (\emph{linear feedback shift register}) de Galois
de 32 bits con el polinomio $x^{32} + x^{22} + x^{2} + x + 1$, que es de período
máximo: recorre los $2^{32}-1$ estados distintos de cero antes de repetirse
\cite{xapp052}. Un paso del LFSR es:

\[
  s \leftarrow
  \begin{cases}
    (s \ll 1) \oplus \texttt{0x00400007} & \text{si el bit 31 de } s \text{ es } 1, \\
    (s \ll 1) & \text{si no.}
  \end{cases}
\]

La constante \code{0x00400007} tiene en 1 los bits 22, 2, 1 y 0, que corresponden a
los términos del polinomio bajo $x^{32}$. El código~\ref{lst:lfsr} implementa el paso
sin desplazamientos, con las dos equivalencias de la sección anterior.

\begin{lstlisting}[language=rvasm, caption={Paso del LFSR y elección del LED en \code{pick\_target}.}, label={lst:lfsr}]
pt_step:
    add  t0, s1, s1                 # t0 = s1 << 1
    bge  s1, zero, pt_no_tap        # bit 31 era 0: sin realimentacion
    xor  t0, t0, t2                 # t2 = LFSR_TAPS
pt_no_tap:
    mv   s1, t0
    addi t1, t1, -1
    bnez t1, pt_step                # LFSR_STEPS = 32 veces

    andi t0, s1, 3                  # indice del LED: 0..3
\end{lstlisting}

\paragraph{Semilla.} Un LFSR es determinista: con la misma semilla produce siempre la
misma secuencia. Si la semilla fuera el contador leído en un punto fijo del programa
tras el reset, sería la misma en cada encendido y el juego repetiría la secuencia de
LEDs. Por eso el juego espera que el jugador apriete un botón para empezar
(\code{wait\_start}), y la semilla es el valor del contador en ese instante. Un humano
no puede repetir un instante con precisión de 40~ns, así que cada partida parte de
un estado distinto. Se le aplica un \code{ori} con 1 para que la semilla nunca sea
cero, el único estado del que un LFSR no sale.

\paragraph{Por qué 32 pasos por ronda.} No todos los bits de la semilla son
aleatorios. Los bits altos del contador cambian lento: el bit 30 vale 0 durante los
primeros 43 segundos tras el encendido. Los bits más bajos tampoco son libres, porque
el bucle que lee el botón lo hace cada una cantidad fija de ciclos. Si el LFSR avanzara
pocos pasos por ronda, los dos bits que eligen el LED vendrían de los bits poco
aleatorios de la semilla. Avanzando 32 pasos, una palabra completa, cada bit del
estado pasa a depender de los bits medios de la semilla, que sí dependen de cuándo
apretó el jugador.

Esto se comprobó con un modelo del LFSR en Python (\code{sim/lfsr\_modelo.py}), que
simula 20.000 partidas con semillas como las que produce el programa: apretones entre
0,5 y 5 segundos tras el encendido, leídos cada $P$ ciclos
(tabla~\ref{tab:lfsr}). Con 2 pasos por ronda, el primer LED es \emph{siempre} el
mismo. Con 32 pasos, cada LED sale un 25\,\% de las veces, los 16 pares posibles de LEDs
consecutivos aparecen con la misma frecuencia, y el número de secuencias distintas
coincide con el que se espera de un generador ideal: con $4^{10} \approx 10^6$
secuencias posibles, en 20.000 partidas se esperan unas 190 repeticiones por azar.

\begin{table}[H]
\centering
\caption{Modelo de 20.000 partidas. $P$: cada cuántos ciclos lee el programa los botones.}
\label{tab:lfsr}
\begin{tabular}{ccccc}
\toprule
Pasos por ronda & $P$ & Primer LED: 0 / 1 / 2 / 3 (\%) & Pares distintos & Secuencias distintas \\
\midrule
2 & 1 & 100 / 0 / 0 / 0 & 1 de 16 & 14\,133 \\
2 & 16 & 100 / 0 / 0 / 0 & 1 de 16 & 14\,218 \\
32 & 1 & 25,5 / 24,9 / 25,0 / 24,6 & 16 de 16 & 19\,809 \\
32 & 16 & 24,7 / 25,2 / 25,1 / 25,0 & 16 de 16 & 19\,796 \\
\bottomrule
\end{tabular}
\end{table}

Cada paso toma 12 o 13 ciclos, según si hay realimentación, así que avanzar 32 pasos
cuesta unos 400 ciclos, 16~$\mu$s. Ocurren antes de encender el objetivo y por lo tanto
no afectan la medición.

\subsection{Instrucciones utilizadas}

La tabla~\ref{tab:instrucciones} lista las instrucciones de máquina que genera el
assembler para \code{game.s}, después de expandir las pseudo-instrucciones. El programa
usa 17 instrucciones distintas de RV32I y ninguna fuera de RV32E.

\begin{table}[H]
\centering
\caption{Instrucciones de máquina de \code{game.hex} (114 en total).}
\label{tab:instrucciones}
\begin{tabularx}{\linewidth}{l c X}
\toprule
Instrucción & Veces & Uso principal \\
\midrule
\code{addi} & 32 & constantes (\code{li}), copias (\code{mv}), contadores \\
\code{jal} & 18 & llamadas (\code{call}) y saltos (\code{j}) \\
\code{lw} & 12 & leer el contador de ciclos y los botones \\
\code{add} & 9 & multiplicar por 2 (máscara, LFSR, decenas) y acumular \\
\code{lui} & 7 & parte alta de constantes grandes (\code{li}) \\
\code{sub} & 6 & intervalos de tiempo y divisiones por restas \\
\code{bltu} & 6 & comparaciones sin signo de tiempos y cocientes \\
\code{jalr} & 6 & retorno de subrutina (\code{ret}) \\
\code{sw} & 6 & escribir LEDs y displays \\
\code{beq}, \code{bne} & 3, 3 & \code{beqz}/\code{bnez} y comparación de botón \\
\code{bge}, \code{bgeu} & 1, 1 & bit 31 del LFSR; saturación a 99 (\code{bleu}) \\
\code{andi}, \code{ori} & 1, 1 & índice del LED; semilla distinta de cero \\
\code{xor}, \code{xori} & 1, 1 & realimentación del LFSR; alternar el parpadeo \\
\bottomrule
\end{tabularx}
\end{table}

\section{El assembler}
\label{sec:assembler}

El assembler es el programa \code{asm/pochoasm.py}, escrito en Python sin bibliotecas
externas (unas 400 líneas, incluidos los comentarios). Recibe un archivo de assembly y escribe un archivo
\code{.hex} con una palabra de 32 bits en hexadecimal por línea. La palabra de la línea
$k$ queda en la dirección $4k$ de la RAM, que es lo que espera \code{\$readmemh} en
\code{pochoco\_ram.v}.

\begin{lstlisting}[language=bash, caption={Uso del assembler.}, label={lst:uso}]
python3 asm/pochoasm.py sw/game.s -o sw/game.hex [-l sw/game.lst] [-D NOMBRE=VALOR]
\end{lstlisting}

La opción \code{-l} escribe además un listado con la dirección, la palabra y la línea de
origen de cada instrucción (anexo~\ref{anx:listado}). La opción \code{-D} reemplaza el
valor de una constante \code{.equ} sin editar el archivo; se usa para simular el juego
con tiempos cortos (sección~\ref{sec:verificacion}).

\subsection{Sintaxis aceptada}

Cada línea tiene la forma
\[
  [\,\textit{etiqueta}\texttt{:}\,]~~[\,\textit{instrucción}~~\textit{operando}\texttt{,}~\textit{operando}\ldots\,]~~[\,\texttt{\#}~\textit{comentario}\,]
\]
y cualquiera de las tres partes puede faltar. Una línea puede tener varias etiquetas.

\begin{itemize}
  \item \textbf{Registros}: \code{x0} a \code{x15} o sus nombres ABI (\code{zero},
        \code{ra}, \code{sp}, \code{gp}, \code{tp}, \code{t0}--\code{t2},
        \code{s0}/\code{fp}, \code{s1}, \code{a0}--\code{a5}). Los registros que
        existen en RV32I pero no en RV32E (\code{x16}--\code{x31}, \code{a6},
        \code{s2}, \code{t3}, \ldots) producen un error explicativo.
  \item \textbf{Números}: decimales, hexadecimales (\code{0x}), binarios (\code{0b}),
        con signo opcional. También se acepta el nombre de una constante \code{.equ},
        con signo menos opcional.
  \item \textbf{Acceso a memoria}: \code{desplazamiento(registro)}, por ejemplo
        \code{lw t0, CYCLES(s0)}. Si el desplazamiento se omite vale 0.
  \item \textbf{Destinos de salto}: siempre una etiqueta.
  \item \textbf{Instrucciones}: las 40 de RV32I: las diez de registro-registro
        (\code{add}, \code{sub}, \code{sll}, \code{slt}, \code{sltu}, \code{xor},
        \code{srl}, \code{sra}, \code{or}, \code{and}), las nueve con inmediato
        (\code{addi}, \code{slti}, \code{sltiu}, \code{xori}, \code{ori}, \code{andi},
        \code{slli}, \code{srli}, \code{srai}), cinco cargas, tres almacenamientos,
        seis branches, \code{lui}, \code{auipc}, \code{jal}, \code{jalr}, \code{ecall},
        \code{ebreak} y \code{fence}.
  \item \textbf{Directivas}: \code{.equ NOMBRE, valor} define una constante. Las
        directivas \code{.text}, \code{.section}, \code{.global} y \code{.globl} se
        aceptan y se ignoran: sirven al assembler de GNU para enlazar varios archivos,
        y aquí hay un solo programa en una sola memoria desde la dirección 0.
\end{itemize}

Las pseudo-instrucciones son abreviaturas estándar de RISC-V \cite{riscv-unpriv} que
no existen en el hardware. El assembler las reescribe como instrucciones reales
(tabla~\ref{tab:pseudo}).

\begin{table}[H]
\centering
\caption{Pseudo-instrucciones aceptadas y su traducción.}
\label{tab:pseudo}
\begin{tabular}{ll}
\toprule
Pseudo-instrucción & Instrucción real \\
\midrule
\code{nop} & \code{addi zero, zero, 0} \\
\code{mv rd, rs} & \code{addi rd, rs, 0} \\
\code{not rd, rs} & \code{xori rd, rs, -1} \\
\code{neg rd, rs} & \code{sub rd, zero, rs} \\
\code{li rd, valor} & \code{addi}, \code{lui} o \code{lui} + \code{addi} (sección~\ref{sec:li}) \\
\code{j etiq} & \code{jal zero, etiq} \\
\code{jal etiq}, \code{call etiq} & \code{jal ra, etiq} \\
\code{jr rs} & \code{jalr zero, 0(rs)} \\
\code{ret} & \code{jalr zero, 0(ra)} \\
\code{beqz rs, etiq} / \code{bnez rs, etiq} & \code{beq} / \code{bne rs, zero, etiq} \\
\code{bgt a, b, etiq} / \code{ble a, b, etiq} & \code{blt} / \code{bge b, a, etiq} \\
\code{bgtu a, b, etiq} / \code{bleu a, b, etiq} & \code{bltu} / \code{bgeu b, a, etiq} \\
\bottomrule
\end{tabular}
\end{table}

En el assembler de GNU, \code{call} se traduce como \code{auipc} + \code{jalr} para
alcanzar cualquier dirección de 32 bits. Aquí se traduce como un solo \code{jal}, cuyo
alcance de $\pm 1$~MB sobra para una RAM de 2~KB.

\subsection{Funcionamiento: tres pasadas}

El problema central de un assembler son las etiquetas usadas antes de definirse: en
\code{call show\_number}, la subrutina está al final del archivo y todavía no se sabe
en qué dirección quedará. Por eso el texto se recorre tres veces:

\begin{enumerate}
  \setcounter{enumi}{-1}
  \item \textbf{Constantes.} Se separa cada línea en etiquetas, mnemónico y operandos, y
        se evalúan las \code{.equ} en orden.
  \item \textbf{Direcciones.} Se recorre el programa con un contador de posición que
        parte en 0. Cada etiqueta recibe el valor actual del contador, y cada
        instrucción lo avanza en 4 bytes. Las pseudo-instrucciones se expanden aquí
        porque determinan cuántas palabras ocupa cada línea: \code{li} puede ocupar una
        o dos.
  \item \textbf{Codificación.} Con todas las direcciones conocidas, cada instrucción
        real se traduce a sus 32 bits.
\end{enumerate}

El tamaño de \code{li} se decide en la primera pasada y depende del valor, así que
\code{li} solo acepta números y constantes \code{.equ}, cuyo valor ya se conoce en ese
momento; no acepta etiquetas.

\subsection{Codificación}

RISC-V tiene seis formatos de instrucción \cite{riscv-unpriv, patterson-riscv}. Todos
dejan el código de operación en los bits 6:0 y los registros siempre en las mismas
posiciones: \code{rd} en 11:7, \code{rs1} en 19:15 y \code{rs2} en 24:20. Así el
hardware puede leer los registros antes de saber qué instrucción es. Lo que cambia es
dónde se reparte el inmediato (tabla~\ref{tab:formatos}).

\begin{table}[H]
\centering
\caption{Formatos de instrucción de RV32I, del bit 31 al 0.}
\label{tab:formatos}
\small
\begin{tabular}{l|c|c|c|c|c|c|}
\cline{2-7}
 & 31 \ldots 25 & 24 \ldots 20 & 19 \ldots 15 & 14 \ldots 12 & 11 \ldots 7 & 6 \ldots 0 \\
\cline{2-7}
R & funct7 & rs2 & rs1 & funct3 & rd & opcode \\
\cline{2-7}
I & \multicolumn{2}{c|}{imm[11:0]} & rs1 & funct3 & rd & opcode \\
\cline{2-7}
S & imm[11:5] & rs2 & rs1 & funct3 & imm[4:0] & opcode \\
\cline{2-7}
B & imm[12\textbar10:5] & rs2 & rs1 & funct3 & imm[4:1\textbar11] & opcode \\
\cline{2-7}
U & \multicolumn{4}{c|}{imm[31:12]} & rd & opcode \\
\cline{2-7}
J & \multicolumn{4}{c|}{imm[20\textbar10:1\textbar11\textbar19:12]} & rd & opcode \\
\cline{2-7}
\end{tabular}
\end{table}

En el assembler cada formato es una función de una línea que coloca cada campo con
desplazamientos y máscaras (código~\ref{lst:enc}). Los formatos B y J no guardan el
bit 0 del desplazamiento, porque las instrucciones están alineadas y siempre es 0. Por
eso duplican el alcance.

\begin{lstlisting}[language=Python, caption={Codificación de los formatos I y B en \code{pochoasm.py}.}, label={lst:enc}]
def enc_i(imm, rs1, f3, rd, op):
    return (imm & 0xFFF) << 20 | rs1 << 15 | f3 << 12 | rd << 7 | op

def enc_b(off, rs2, rs1, f3, op):
    # imm[12|10:5] en 31:25, imm[4:1|11] en 11:7 (el bit 0 no se guarda: siempre es 0)
    return (((off >> 12) & 1) << 31 | ((off >> 5) & 0x3F) << 25 | rs2 << 20 | rs1 << 15
            | f3 << 12 | ((off >> 1) & 0xF) << 8 | ((off >> 11) & 1) << 7 | op)
\end{lstlisting}

La información de cada instrucción está en una tabla de datos, no en código: el
mnemónico apunta a su formato, \code{opcode}, \code{funct3} y \code{funct7}. Por ejemplo,
\code{'bne': ('B', 0x63, 1, 0)}. Agregar una instrucción es agregar una fila.

\subsection{Etiquetas e inmediatos}

\paragraph{Etiquetas.} Los branches y \code{jal} no guardan la dirección de destino
sino la distancia desde la instrucción que salta: $\text{desplazamiento} =
\text{dirección de la etiqueta} - pc$. El procesador suma ese valor a su \code{pc}. Un
desplazamiento negativo es un salto hacia atrás, como el de un bucle.

\paragraph{Rangos.} Cada inmediato se verifica contra lo que cabe en su campo, con
signo: $[-2048, 2047]$ para los formatos I y S, $[0, 31]$ para la cantidad de
desplazamiento, $[0, 2^{20}-1]$ para \code{lui} y \code{auipc}, $[-4096, 4094]$ para
los branches y $[-2^{20}, 2^{20}-2]$ para \code{jal}.

\paragraph{La pseudo-instrucción \code{li}.}
\label{sec:li}
Carga un valor de 32 bits en un registro. Si cabe en 12 bits con signo, basta un
\code{addi} desde \code{zero}. Si no, se parte en una parte alta de 20 bits, que carga
\code{lui}, y una baja de 12 bits, que suma \code{addi}. La dificultad es que
\code{addi} interpreta sus 12 bits \emph{con signo}. Si el bit 11 de la parte baja
es 1, \code{addi} resta en vez de sumar, y la parte alta tiene que compensarlo sumando 1.
El assembler calcula:

\[
  \text{lo} = \text{extensión de signo de } (v \mathbin{\&} \texttt{0xFFF}), \qquad
  \text{hi} = \frac{v - \text{lo}}{2^{12}} \bmod 2^{20}.
\]

Si $\text{lo} = 0$ basta el \code{lui}. Este caso se verificó ejecutando en Python el
par \code{lui}/\code{addi} generado para 2.017 valores, incluidos los bordes
(\code{0x7FF}, \code{0x800}, \code{0xFFF}, \code{0x80000000}, \code{0xFFFFFFFF}) y
2.000 al azar.

\subsection{Validación}

Cada error se reporta con el archivo, la línea y el texto de la línea, y el assembler
sigue revisando el resto del archivo para mostrar todos los errores de una vez. Si hay
algún error no escribe el \code{.hex}, para que nunca quede en disco un archivo a medio
generar. El código~\ref{lst:errores} muestra un ejemplo.

\begin{lstlisting}[caption={Salida del assembler ante un programa con errores.}, label={lst:errores}]
malo.s:3: error: 'x16' no existe en RV32E: solo hay 16 registros (x0-x15)
    add  x16, a0, a1
malo.s:7: error: inmediato de 12 bits fuera de rango: 3000 (permitido -2048 a 2047)
    addi t0, t0, 3000
malo.s:9: error: etiqueta no definida: 'fin'
    j    fin
3 error(es); no se escribio ningun archivo.
\end{lstlisting}

Los errores detectados son: registro inexistente o fuera de RV32E, instrucción
desconocida, cantidad incorrecta de operandos, operando vacío, inmediato o salto fuera
de rango, acceso a memoria mal escrito, constante no definida, etiqueta no definida o
repetida, constante repetida, nombre de constante reservado, directiva no soportada y
programa más grande que la RAM de 512 palabras.

Además, el assembler \emph{avisa} (sin fallar) si el programa usa un desplazamiento,
porque el Espino Core lo ejecutaría como una suma y el resultado sería distinto del que
espera quien lee el código.

\subsection{Ejemplo de traducción}

\paragraph{Un branch.} La instrucción \code{bne t0, a3, wrong} de \code{game.s} está en
la dirección \code{0x5C}, y la etiqueta \code{wrong} quedó en \code{0xA8}.

\begin{enumerate}
  \item Desplazamiento: $\texttt{0xA8} - \texttt{0x5C} = \texttt{0x4C} = 76$. En 13
        bits: \code{0\,0000\,0100\,1100}.
  \item Campos: \code{bne} es formato B con opcode \code{1100011} y funct3 \code{001};
        \code{t0} es \code{x5} (\code{00101}) y \code{a3} es \code{x13} (\code{01101}).
  \item Se reparte el desplazamiento: imm[12] = 0, imm[10:5] = \code{000010},
        imm[4:1] = \code{0110}, imm[11] = 0.
\end{enumerate}

\begin{table}[H]
\centering
\small
\begin{tabular}{|c|c|c|c|c|c|c|c|}
\hline
imm[12] & imm[10:5] & rs2 & rs1 & funct3 & imm[4:1] & imm[11] & opcode \\
\hline
0 & 000010 & 01101 & 00101 & 001 & 0110 & 0 & 1100011 \\
\hline
\end{tabular}
\end{table}

Agrupando de a cuatro bits:
\code{0000 0100 1101 0010 1001 0110 0110 0011} = \code{0x04D29663}, que es la línea 24
de \code{game.hex}.

\paragraph{Una constante grande.} \code{li a0, WAIT\_CYCLES}, con
$\text{WAIT\_CYCLES} = 75\,000\,000 = \texttt{0x047868C0}$:

\begin{enumerate}
  \item Los 12 bits bajos son \code{0x8C0}. Su bit 11 es 1, así que con signo valen
        $\texttt{0x8C0} - 4096 = -1856$.
  \item La parte alta se redondea hacia arriba para compensar:
        $(75\,000\,000 + 1856) / 4096 = 18\,311 = \texttt{0x4787}$.
  \item \code{lui a0, 0x4787} carga $\texttt{0x04787000} = 75\,001\,856$, y
        \code{addi a0, a0, -1856} deja $75\,000\,000$.
\end{enumerate}

Codificadas: \code{lui} (formato U, opcode \code{0110111}, rd = \code{01010}) da
\code{0x04787537}, y \code{addi} (formato I, inmediato \code{0x8C0}, rs1 = rd =
\code{01010}, opcode \code{0010011}) da \code{0x8C050513}. Son las líneas 13 y 14 de
\code{game.hex}.

\section{Verificación}
\label{sec:verificacion}

La verificación tiene cuatro niveles, de lo más aislado a lo más completo. Todas las
pruebas son automáticas y se ejecutan con \code{make test}, \code{make sim} y
\code{make sim-post}.

\subsection{Pruebas del assembler}

El archivo \code{asm/test\_pochoasm.py} contiene 20 pruebas, escritas con el módulo
\code{unittest} de Python (tabla~\ref{tab:pruebas-asm}). Las dos primeras filas son las
más importantes, porque comparan contra algo que no escribió el grupo.

\begin{itemize}
  \item \textbf{Referencia GNU.} El repositorio original trae tres programas de ejemplo
        (\code{blink.s}, \code{7seg.s}, \code{buttons\_leds.s}) con sus \code{.hex}
        generados por el assembler oficial de RISC-V de GNU. Se guardaron en
        \code{asm/golden/}, y el assembler del grupo produce exactamente las mismas 31
        palabras.
  \item \textbf{Ida y vuelta contra el decodificador.} Para cada formato con inmediato
        se codifica \emph{cada} valor posible y se decodifica con las fórmulas copiadas
        de \code{espino\_decoder.v} (por ejemplo,
        \code{imm\_b = \{\{19\{i[31]\}\}, i[31], i[7], i[30:25], i[11:8], 1'b0\}}). Tiene
        que volver el mismo número: 4.096 valores en I y S, 4.096 desplazamientos en B y
        1.048.576 en J. Si el assembler pusiera un bit en la posición equivocada, el
        procesador leería otro valor y esta prueba lo detectaría.
\end{itemize}

\begin{table}[H]
\centering
\caption{Pruebas del assembler (\code{make test}).}
\label{tab:pruebas-asm}
\begin{tabularx}{\linewidth}{l c X}
\toprule
Grupo & Pruebas & Qué se comprueba \\
\midrule
Referencia GNU & 1 & los 3 ejemplos del repositorio, palabra por palabra \\
Ida y vuelta & 5 & todos los inmediatos de I, S, B y J contra el decodificador; saltos hacia adelante y hacia atrás con etiquetas \\
Palabras conocidas & 9 & \code{nop}, \code{ret}, \code{ecall} y otras palabras de la especificación; todas las pseudo-instrucciones; \code{li} con 2.017 valores; \code{.equ}; \code{-D}; aviso de desplazamientos; comentarios y etiquetas \\
Validación & 4 & 21 programas erróneos, cada uno con su mensaje y su línea; reporte de todos los errores juntos \\
El juego & 1 & \code{game.s} se ensambla sin avisos y cabe en la RAM \\
\bottomrule
\end{tabularx}
\end{table}

\subsection{Prueba de instrucciones sobre el procesador}

El programa \code{sim/isa\_test.s} ejercita cada instrucción de RV32I en 22 pruebas:
operaciones de ALU con valores con y sin signo, inmediatos en los extremos de su rango,
\code{lui} y \code{auipc}, cada branch tomado y no tomado, \code{jal} y \code{jalr},
cargas y almacenamientos de palabras, medias palabras y bytes con extensión de signo,
lectura de los botones y del contador de ciclos, el registro \code{x0} y la ejecución de
un desplazamiento como suma. Cada prueba compara su resultado con el esperado y, si
falla, escribe su número en los displays.

El programa se ensambla con el assembler del grupo y se ejecuta sobre el RTL completo
del SoC en Icarus Verilog (\code{sim/isa\_tb.v}), usando los modelos de simulación de
Lattice para los bloques de RAM. Así se prueba el assembler contra el hardware real: si
una instrucción se codificara mal, el procesador ejecutaría otra cosa. Resultado:
\textbf{las 22 pruebas pasan}.

Para confirmar que la prueba detecta errores de verdad, se introdujeron dos errores a
propósito en el assembler. Al cambiar el \code{funct3} de \code{lbu} por el de
\code{lb}, falla la prueba 15. Al perder un bit del inmediato del formato S, el
programa se escribe encima de sí mismo y no termina. En ambos casos la simulación
reporta el fallo.

\subsection{Simulación del juego}

El testbench \code{sim/game\_tb.v} juega dos partidas completas contra el SoC corriendo
\code{game.s}, y comprueba cada regla del enunciado mirando solo los pines de la placa:
los cuatro LEDs y los catorce segmentos de los displays, que decodifica de vuelta a
dígitos igual que los leería una persona. Aprieta y suelta los botones con rebotes: tres
pulsos cortos antes de quedar firmes.

Simular 3 segundos reales son 75 millones de ciclos, demasiado para un simulador. Por
eso el testbench usa el mismo \code{game.s} ensamblado con constantes más cortas,
mediante la opción \code{-D} del assembler: \code{WAIT\_CYCLES}~=~20.000,
\code{TENTH\_CYCLES}~=~1.000, \code{DEBOUNCE\_CYCLES}~=~200 y
\code{BLINK\_CYCLES}~=~2.000. El programa es idéntico; solo cambian esos cuatro números.

La primera partida usa diez tiempos de reacción distintos, elegidos para ejercitar los
casos borde: un dígito, dos dígitos, y 120 décimas, que no cabe en los displays. Además,
en la ronda 3 aprieta primero un botón incorrecto, y en la ronda 7 mantiene un botón
apretado durante toda la fase de los cuatro LEDs. La tabla~\ref{tab:sim-juego} resume lo
que se comprobó.

\begin{table}[H]
\centering
\caption{Revisiones del testbench del juego (\code{make sim}). Todas pasan.}
\label{tab:sim-juego}
\begin{tabularx}{\linewidth}{X l}
\toprule
Regla & Resultado observado \\
\midrule
Al encender: displays en \code{00} y LEDs parpadeando & \code{00}, 0101 $\leftrightarrow$ 1010 \\
Los cuatro LEDs duran al menos \code{WAIT\_CYCLES} en cada ronda & sí, en las 21 rondas \\
Luego se enciende un solo LED & sí \\
Tiempo mostrado: 3, 7, 12, 5, 120, 1, 15, 4, 8, 6 décimas & \code{03 07 12 05 99 01 15 04 08 06} \\
Botón incorrecto: \code{EE} y la ronda se repite & \code{EE}, vuelven los 4 LEDs \\
El error no cuenta como ronda ni suma al promedio & la partida termina tras 10 correctas \\
Botón sostenido durante los 4 LEDs: no es error y el objetivo espera a que se suelte & sí, con 200 ciclos sin botones \\
Tras la décima: el último tiempo sigue visible \code{WAIT\_CYCLES} & \code{06} con los LEDs apagados \\
Promedio: $\lfloor 181 / 10 \rfloor = 18$ décimas & \code{18}, LEDs parpadeando \\
Una partida nueva empieza en \code{00} & sí \\
Dos partidas empezadas en instantes distintos dan secuencias distintas & \code{2 2 3 0 1 2 2 2 3 0} vs \code{0 2 1 1 2 3 0 3 3 0} \\
\bottomrule
\end{tabularx}
\end{table}

\subsection{Simulación post-síntesis}

Las simulaciones anteriores corren sobre el RTL, la descripción en Verilog. Lo que se
graba en la FPGA es otra cosa: el circuito que produce Yosys, hecho de tablas de
búsqueda, flip-flops y bloques de RAM de la iCE40, con el programa ya cargado en esos
bloques. Si la síntesis interpretara algo distinto del simulador, por ejemplo la
escritura por bytes de la RAM o su inicialización con \code{\$readmemh}, la
simulación del RTL no lo mostraría.

\code{make sim-post} sintetiza el diseño con la versión corta del juego, exporta el
circuito resultante como Verilog de celdas y le aplica el mismo testbench. El resultado
es idéntico al de la tabla~\ref{tab:sim-juego}: los mismos tiempos mostrados, el mismo
\code{EE}, el mismo promedio \code{18} y las mismas secuencias de LEDs en ambas
partidas. La simulación toma unos 7 minutos.

\subsection{Problemas encontrados}

El testbench encontró dos problemas durante el desarrollo, ambos en el propio
testbench y no en el juego. El primero: medía la duración de la fase de cuatro LEDs
desde que empezaba a mirarla, y no desde que los LEDs cambiaban, así que reportaba
fases cortas que no lo eran. El segundo: cambiaba los botones en el mismo instante del
flanco de reloj en que el SoC los lee, una carrera que hacía variar la semilla entre una
compilación y otra. Se corrigió asignando los botones con asignación no bloqueante.

\section{Compilación, programación y cambio de constantes}
\label{sec:operacion}

Todas las herramientas vienen en OSS CAD Suite \cite{osscad}: Yosys, nextpnr, IceStorm,
Icarus Verilog y Python 3. Desde su terminal, en la carpeta del proyecto:

\begin{table}[H]
\centering
\caption{Comandos del Makefile.}
\label{tab:make}
\begin{tabularx}{\linewidth}{l X}
\toprule
Comando & Qué hace \\
\midrule
\code{make prog} & ensambla \code{game.s}, sintetiza, hace place and route, genera \code{game\_top.bin} y programa la Go Board \\
\code{make} & lo mismo, sin programar \\
\code{make hex} & solo \code{sw/game.s} $\to$ \code{sw/game.hex} \\
\code{make test} & pruebas del assembler \\
\code{make sim} & prueba de instrucciones y simulación del juego \\
\bottomrule
\end{tabularx}
\end{table}

El Makefile conoce las dependencias: \code{game\_top.bin} depende de \code{game.hex},
que depende de \code{game.s} y del assembler. Por eso, tras editar \code{game.s}, basta
\code{make prog}: regenera solo lo necesario. La compilación completa toma unos
8 segundos.

Para regenerar \code{game.hex} sin make:

\begin{lstlisting}[language=bash]
python3 asm/pochoasm.py sw/game.s -o sw/game.hex
\end{lstlisting}

\paragraph{Cambiar una constante.} Durante la demostración se pedirá modificar una
constante del juego. Todas están con nombre al comienzo de \code{sw/game.s}
(código~\ref{lst:constantes}). Por ejemplo, para jugar cinco rondas en vez de diez:

\begin{enumerate}
  \item Editar la línea \code{.equ ROUNDS, 10} a \code{.equ ROUNDS, 5}.
  \item Ejecutar \code{make prog}.
\end{enumerate}

El promedio se sigue calculando bien porque divide por \code{ROUNDS}, no por 10. De la
misma forma, \code{.equ WAIT\_CYCLES, 125000000} deja la fase de los cuatro LEDs en
5 segundos. El \code{.hex} cambia porque la RAM se inicializa al sintetizar, así que un
cambio en el programa exige volver a sintetizar el diseño completo.

\section{Conclusiones}

Se implementó el juego de reflejos completo sobre el Pochoco SoC: temporización de tres
segundos, elección pseudoaleatoria del LED objetivo, medición del tiempo de respuesta
en décimas de segundo, manejo de botones incorrectos y promedio de diez rondas. Todo
el comportamiento se verificó en simulación sobre el RTL del SoC y sobre el circuito
sintetizado, el diseño sintetiza para la iCE40 HX1K con margen de frecuencia, y el juego
funciona en la Go Board.

El único cambio de hardware necesario fue un contador de ciclos de 32 bits. Con él, toda
medición de tiempo se reduce a una resta sin signo, que además tolera que el contador
dé la vuelta. La alternativa, bucles de espera calibrados, habría acoplado la duración
de cada fase del juego al número exacto de instrucciones del programa.

Las restricciones del procesador obligaron a expresar operaciones elementales con las
instrucciones disponibles: multiplicar por dos es sumar un número consigo mismo, leer
el bit 31 es preguntar si el número es negativo y dividir es restar repetidamente. El
mismo criterio permitió implementar un LFSR sin desplazamientos. El modelo del
generador mostró que el número de pasos entre elecciones importa: con pocos pasos el
primer LED de cada partida era siempre el mismo, aunque la semilla viniera del
contador.

La prueba en la placa mostró el límite de la simulación. El juego pasaba todas las
simulaciones, incluida la del circuito sintetizado, y aun así no arrancaba en la FPGA,
porque el reset de encendido original soltaba al procesador antes de que la RAM
estuviera disponible. Los modelos de simulación no representan ese retraso. La causa
se aisló grabando programas de prueba cada vez más cercanos al juego y descartando una
hipótesis por vez, y la corrección fue una línea de Verilog.

En el assembler, el problema de fondo son las referencias hacia adelante: una etiqueta
se usa antes de saber su dirección. Se resolvió con pasadas separadas para las
direcciones y para la codificación. La decisión que más confianza dio sobre su
corrección fue probarlo contra referencias externas: la salida del assembler de GNU,
las fórmulas del decodificador del procesador y el procesador mismo, en simulación.


\newpage
\section*{Referencias}
\addcontentsline{toc}{section}{Referencias}
\printbibliography[heading=none]

\newpage
\appendix
\section{Programa completo: \code{sw/game.s}}
\label{anx:game}

\lstinputlisting[language=rvasm]{anexos/game.s}

\newpage
\section{Listado de traducción: \code{sw/game.lst}}
\label{anx:listado}

Generado con \code{python3 asm/pochoasm.py sw/game.s -o sw/game.hex -l sw/game.lst}.
Cada línea tiene la dirección, la palabra de 32 bits que queda en \code{game.hex} y la
línea de \code{game.s} que la produjo. Las pseudo-instrucciones que se expanden a dos
palabras (\code{li} de valores grandes) muestran el texto solo en la primera.

{\scriptsize
\lstinputlisting[language=rvasm, basicstyle=\ttfamily\scriptsize]{anexos/game.lst}
}


\newpage
\section*{Declaración de Uso de Herramientas de Inteligencia Artificial}
\addcontentsline{toc}{section}{Declaración de Uso de Herramientas de Inteligencia Artificial}

% ============================================================================
% REVISAR ANTES DE FIRMAR. Esta es una declaracion formal y debe reflejar
% exactamente lo que hizo el grupo. Ajustar el alcance, los integrantes que
% usaron la herramienta y las fechas antes de la entrega.
% ============================================================================

\subsection*{1. Identificación de la actividad}
\begin{itemize}
    \item \textbf{Curso:} Arquitectura de Computadores
    \item \textbf{Código del curso:} 433
    \item \textbf{Nombre de la actividad:} Proyecto 2: Juego en Assembly sobre Pochoco SoC
    \item \textbf{Fecha de entrega:} Martes 29 de septiembre de 2026
    \item \textbf{Nombre(s) de(los) estudiante(s):} Josefa Alcaíno Lira, Felipe Izcúe Ferrer, Martín Norton Haigis
\end{itemize}

\subsection*{2. Declaración sobre uso de IA}
\noindent $\square$ No se utilizaron herramientas de inteligencia artificial en ninguna etapa del desarrollo de esta actividad.

\noindent $\boxtimes$ Sí se utilizaron herramientas de inteligencia artificial, según se detalla en la sección 3 de este documento.

\subsection*{3. Descripción de uso de IA}
\begin{table}[H]
    \centering
    \begin{tabularx}{\linewidth}{|l|X|}
        \hline
        \textbf{Herramienta empleada} & Claude (Anthropic), a través de la interfaz Claude Code. \\
        \hline
        \textbf{Fechas de uso} & 27 y 28 de septiembre de 2026. \\
        \hline
        \textbf{Estudiantes} & Martín Norton Haigis. \\
        \hline
        \textbf{Propósito del uso} &
        Desarrollo, con la herramienta, del código del proyecto a partir del enunciado
        y del repositorio base: el contador de ciclos y el módulo \code{game\_top} en
        Verilog, el programa del juego en assembly, el assembler en Python, sus pruebas
        automáticas, los testbenches de simulación y el Makefile.
        Diseño y ejecución de la verificación: simulación en Icarus Verilog, medición
        de ciclos por instrucción, modelo del generador pseudoaleatorio y síntesis con
        Yosys y nextpnr.
        Redacción de este informe y de la documentación del repositorio. \\
        \hline
        \textbf{Descripción de integración} &
        Las decisiones de diseño se tomaron durante el desarrollo y quedaron
        registradas en el código y en este informe. Todo el código se verificó con las
        pruebas y simulaciones descritas en la sección~\ref{sec:verificacion}, cuyos
        resultados se reportan tal como se obtuvieron. El grupo revisó el código y el
        texto, y comprende y puede explicar cada parte del diseño, del programa y del
        assembler. \\
        \hline
    \end{tabularx}
\end{table}

\subsection*{4. Declaración de responsabilidad académica}
\noindent Los firmantes de este documento declaramos que:
\begin{itemize}
    \item Hemos revisado críticamente todo el contenido entregado.
    \item Comprendemos plenamente el trabajo presentado y somos capaces de explicarlo y defenderlo académicamente.
    \item Asumimos responsabilidad por la calidad, veracidad y originalidad del contenido entregado, independientemente del uso o no uso de herramientas de inteligencia artificial.
    \item La información declarada en este documento es veraz.
\end{itemize}

\vspace{1.5cm}
\noindent
\begin{tabular}{p{4.6cm}p{4.6cm}p{4.6cm}}
    \vspace{2cm} & \vspace{2cm} & \vspace{2cm} \\
    \hrulefill & \hrulefill & \hrulefill \\
    Josefa Alcaíno Lira & Felipe Izcúe Ferrer & Martín Norton Haigis \\
\end{tabular}


\end{document}
